\begin{helper}[At most three retained departures]
For a defective color $\delta$, put
$X=V(T)$ and
$Z=X\setminus\{x^0_{\gamma(\delta)},x^1_{\gamma(\delta)}\}$, and let
$B$ be the finite set of exact $\delta$-colored boundary-edge identities
of $Z$ which have an incident vertex outside $X$.  Then $|B|\le3$.
\end{helper}
\begin{proof}
Fix $\delta$, put $X=V(T)$, $i=\gammaIndex(\delta)$,
$\gamma=\gamma_i$, and use the sets $Z,B$ of the statement.  The
codomain clause for $\gammaIndex$ in \noderef{ScheideDefectiveSetup}
gives $i<s$.  By
\noderef{ScheideReplacementBoundaryParity}, $|B|$ is odd,
$f^0_\delta\in B$, and $v^0_\delta\in Z$.  Suppose $|B|>3$.  Oddness
gives $|B|\ge5$, so there are three pairwise distinct exact identities
$g_1,g_2,g_3$ in $B\setminus\{f^0_\delta\}$.

Membership in $B$, expanded through \noderef{BoundaryEdges},
\noderef{Incident}, and \noderef{LooplessMultigraph}, orients each edge
as
\[
 \operatorname{ends}(g_j)=\{z_j,w_j\},\qquad
 z_j\in Z,\quad w_j\notin X,\quad\varphi(g_j)=\delta.
\]
The vertices $w_1,w_2,w_3$ are pairwise distinct: equality of two would
make the corresponding distinct $\delta$-identities adjacent, contrary
to \noderef{AdjacentEdges} and \noderef{IsProperEdgeColoring}; this
properness is supplied by the global field of
\noderef{ScheideReducedConfiguration}.

Let $\varphi_\delta$ be the $(\gamma,\delta)$ component switch rooted at
$v^0_\delta$.  The exact clauses of
\noderef{ScheideDefectiveSwitchEffect} say that it is proper, leaves
exactly $e_0$ uncolored, fixes every identity incident with $X$ other
than $f^0_\delta$, and makes $\delta$ missing at $v^0_\delta$.
Consequently every $g_j$ retains color $\delta$.

Put $n=T.\operatorname{length}$ and
$a=p-1+2i$, the zero-based position of the first selected block edge.
Using $i<s$ from \noderef{ScheideDefectiveSetup}, together with $p\ge2$
and the length and position equations in
\noderef{ScheideStructuralConfiguration}, gives $a\ge1$, proves that
$a,a+1$ are valid edge positions and $a+1,a+2$ are valid vertex
positions, identifies the former with the two selected block edges, and
identifies the latter with the two selected block vertices.  Define an
ordered sequence $T^+$ by retaining all entries strictly before these
positions, shifting every later old block pair two positions to the
left, and appending
$(g_1,w_1),(g_2,w_2),(g_3,w_3)$.

We verify \noderef{TreeSequence} exactly.  Its length is $n+1$ and its
first edge is the old first edge because $a\ge1$.  Old retained vertices
and edges remain injective.  By \noderef{TreeSequenceVertexSet}, $X$ is
exactly the image of the old vertex map.  Hence $w_j\notin X$ makes each
new vertex fresh from every old displayed vertex, and the new vertices
are pairwise distinct.  The attachment equation in
\noderef{TreeSequence} puts both endpoints of every old displayed edge
in that vertex-map image $X$, so none of $g_1,g_2,g_3$ is an old edge;
the chosen identities are also pairwise distinct.  Attachments before
$a$ retain their old witnesses.  Every shifted old edge is a later block
edge, and \noderef{ScheideStructuralConfiguration} attaches it directly
to a retained base vertex.  Finally $g_j$ attaches $w_j$ to the retained
vertex $z_j$.  Thus all tree-sequence fields hold.

This sequence is Tashkinov for $\varphi_\delta$.  For an old retained
edge before $a$, its earlier witness is retained.  Its color is not
$\gamma$: nonfirst base edges avoid every block color and an earlier
block edge has the distinct color $\gamma_r$ with $r<i$.  A shifted
later block edge has color $\gamma_r$ with $r>i$ and retains its base
missing witness.  Hence the switch-effect equalities preserve all these
old witnesses.  Each new edge has color $\delta$, now missing at the
retained earlier vertex $v^0_\delta$.  This proves the
\noderef{TashkinovTree} condition.

The graph-only fields of the old \noderef{KTriple} are unchanged, while
properness and exact-one-uncoloredness come from
\noderef{ScheideDefectiveSwitchEffect}; hence the switched data form a
new $k$-triple.  By \noderef{TreeSequenceVertexSet}, both $V(T)$ and
$V(T^+)$ are the images of their respective vertex maps.  The two
selected block vertices are distinct old image elements, the three
appended vertices are pairwise distinct and fresh from the old image,
and all other old image elements are retained.  Therefore
\[
 |V(T^+)|=|X|-2+3=|X|+1=\noderef{TashkinovOrder}(G)+1,
\]
where the last equality follows from the global maximum field in
\noderef{GlobalMaximumTriple}.  This contradicts the defining upper
bound in \noderef{TashkinovOrder} for a Tashkinov tree of a $k$-triple.
Therefore $|B|\le3$.
\end{proof}
